(1) 对横向小扰动，理想MHD线性化为
\[
\rho\partial_t\boldsymbol u=(B_0/\mu)\partial_z\boldsymbol b,\qquad
\partial_t\boldsymbol b=B_0\partial_z\boldsymbol u.
\]
消去 \(\boldsymbol b\) 得 \(\partial_t^2\boldsymbol u=v_A^2\partial_z^2\boldsymbol u\)，
\[
\boxed{v_A=B_0/\sqrt{\mu\rho},\qquad\omega^2=k^2v_A^2}.
\]
对于 \(e^{i(kz-\omega t)}\) 且 \(\omega=+kv_A\)，感应方程给 \(\boldsymbol b=-\sqrt{\mu\rho}\boldsymbol u\)，符号取决于传播方向。

(2) 定义 \(\boldsymbol z^\pm=\boldsymbol u\pm\boldsymbol b/\sqrt{\mu\rho}\)，则
\[
\partial_t\boldsymbol z^\pm=\pm v_A\partial_z\boldsymbol z^\pm.
\]
因此 \(\boldsymbol z^-\) 满足 \(\partial_t\boldsymbol z^-+v_A\partial_z\boldsymbol z^-=0\)，沿 \(+z\) 传播；\(\boldsymbol z^+\) 沿 \(-z\) 传播。该符号对应来自题设 \(B_0>0\) 与 \(e^{i(kz-\omega t)}\) 约定，不能只按变量上标猜测方向。

(3) 两方程分别点乘 \(\boldsymbol u,\boldsymbol b/\mu\) 相加，得
\[
\boxed{\partial_t\left(\frac{\rho u^2}{2}+\frac{b^2}{2\mu}\right)
+\partial_z\left(-\frac{B_0}{\mu}\boldsymbol u\cdot\boldsymbol b\right)=0}.
\]
正向波 \(\boldsymbol b=-\sqrt{\mu\rho}\boldsymbol u\) 的能流为
\(-B_0\boldsymbol u\cdot\boldsymbol b/\mu=\rho v_Au^2>0\)，同时磁能与动能相等。若误将感应方程符号反写，能流方向会与传播方向冲突。

(4) 定义波阻抗 \(Z_j=\rho_jv_{Aj}=B_0\sqrt{\rho_j/\mu}\)。在界面处横向速度连续 \(u_i+u_r=u_t\)，横向磁场连续给 \(Z_1(u_i-u_r)=Z_2u_t\)。因此
\[
\boxed{r_u=\frac{u_r}{u_i}=\frac{Z_1-Z_2}{Z_1+Z_2},\quad
t_u=\frac{u_t}{u_i}=\frac{2Z_1}{Z_1+Z_2}}.
\]
能流反射率 \(R=r_u^2\)，透射率 \(T=(Z_2/Z_1)t_u^2\)，直接化简 \(R+T=1\)。该结论限定于正入射、理想无耗散、界面 \(B_0,\mu\) 连续且仅有剪切Alfvén分支。
